\begin{center}\begin{tikzpicture}[font=\small,>=stealth,line join=round]
\draw[gray!40,rounded corners=3pt] (0,0) rectangle (14.2,8.3);\node at (7.1,7.700000000000001) {证明三：整条曲线下移3，中心高度变成0};
\begin{scope}[shift={(2.207692307692308,3.102631578947368)},x=0.9230769230769231cm,y=0.4842105263157894cm]\draw[->,gray] (-1.2,0)--(5.3,0) node[right]{$x$};\draw[->,gray] (0,-3)--(0,6.5) node[above]{$y$};\draw[blue,thick,solid,domain=-0.8:4.8,samples=180] plot(\x,{3+2*cos(deg(pi*\x/2))});\draw[orange!85!black,thick,dashed,domain=-0.8:4.8,samples=180] plot(\x,{2*cos(deg(pi*\x/2))});\fill[orange!85!black] (1,3) circle(2.2pt);\fill[orange!85!black] (1,0) circle(2.2pt);\draw[gray,dashed] (-1,3)--(5,3);\draw[<->,violet] (1.2,0)--(1.2,3);\end{scope}\node[anchor=west,align=left,text width=5.0cm] at (8.4,4.0) {蓝：$f=3+2\cos\frac{\pi x}2$\\[9pt]橙虚线：$g=f-3$\\[9pt]中心 $(1,3)\to(1,0)$\\[9pt]周期都为4\\[9pt]纵向平移不改重复距离};
\node[text width=12.6cm,align=center] at (7.1,.65) {把中线当作新的零线，关于中心的高度翻转就变成正负号交换。};
\end{tikzpicture}\end{center}
